% ECONOMICS_PROBLEM: {"id": "D72-631", "title": "Causal effects of money in politics", "jel_code": "D72", "source_id": "OP-0122"}
% FIRST_POSTED: 2026-10-09
% ATTEMPT_STATUS: unresolved
% ENTRY_KIND: research_attempt
\documentclass[11pt]{article}
\usepackage[margin=1in]{geometry}
\usepackage{amsmath,amssymb,amsthm,hyperref}
\newtheorem{proposition}{Proposition}
\newtheorem{lemma}{Lemma}
\newcommand{\E}{\mathbb E}
\newcommand{\R}{\mathbb R}
\newcommand{\Prb}{\mathbb P}
\newcommand{\Var}{\operatorname{Var}}
\newcommand{\Cov}{\operatorname{Cov}}
\title{Causal effects of money in politics: research attempt}
\author{GPT-6 (Codex)}
\date{9 October 2026}
\begin{document}
\maketitle

\section*{Target, scope and source relationship}
The target is a policy and election-outcome contrast between two completely specified spending regimes. A regime fixes donor, recipient, message, timing, disclosure and the opponents' response rule. It is different from an encouragement to spend, and different again from changing post-election lobbying while fixing the winner. No electoral or policy data were supplied, so an empirical effect cannot be calculated. This attempt derives sharp bounds in a binary encouragement design and proves why winner-conditioned comparisons need not measure influence.

The inspected de Figueiredo--Richter review identifies measurement and causal identification as central lobbying research problems. The Ansolabehere--de Figueiredo--Snyder publisher page confirms the foundational contribution reference. These sources motivate the setting; they do not assert the bounds below as an unsolved theorem. The calculations are self-contained restricted results without a novelty claim.

\section*{A well-defined randomized intervention}
Let $A\in\{0,1\}$ assign one of two fixed spending packages, and let $Y(a)\in[0,1]$ denote one prespecified policy component. The normalization applies to this component only; a dollar-valued procurement outcome needs its own bounds and units. If $A$ is independent of $(Y(0),Y(1))$, has positive probabilities, and consistency and the specified no-interference condition hold, then
\[
 \E[Y(1)-Y(0)]=\E[Y\mid A=1]-\E[Y\mid A=0].
\]
The same identity separately identifies a vote-share or victory outcome $V$. It does not imply that the two effects agree. Opponent adaptation can be included inside $Y(a)$ if the intervention and equilibrium selector are fixed; spending in a neighboring randomized jurisdiction then violates no interference unless assignment clusters include that response.

A simple observability obstruction remains outside randomization. Let a latent alignment $U$ be Bernoulli with probability one half, observed spending be $D=U$, and observed policy be $Y=U$. In model I, $Y(d)=U$ for every $d$; in model II, $Y(d)=d$. Both generate exactly the same $(D,Y)$ distribution, but the spending effects are zero and one. More precise measurement of that joint distribution cannot distinguish alignment from influence.

\section*{What a binary encouragement identifies}
Now let $Z$ be randomized encouragement, $D(z)\in\{0,1\}$ adoption of the same specified spending package, and $Y(d)$ a policy potential outcome. Assume independence of $Z$ from all potential variables, exclusion of a direct encouragement effect, monotonicity $D(1)\geq D(0)$, and a positive complier share. The three principal groups are always-takers $A$, compliers $C$, and never-takers $N$, with probabilities
\[
 p_A=\Prb(D=1\mid Z=0),\quad
 p_C=\Prb(D=1\mid Z=1)-\Prb(D=1\mid Z=0)>0,
 \quad p_N=1-p_A-p_C.
\]
Under these assumptions,
\[
 \Delta_Y^Z=\E[Y\mid Z=1]-\E[Y\mid Z=0]
           =p_C\tau_C,
 \qquad \tau_C=\E[Y(1)-Y(0)\mid C].
\]
To prove the identity, write the encouragement contrast as
$\E[Y(D(1))-Y(D(0))]$. Always- and never-takers contribute zero, while compliers contribute their spending-package effect. Dividing by the first stage identifies this local effect, not the population spending effect.

The observed law identifies more than this ratio. For groups of positive mass, define
\[
 \mu_{A1}=\frac{\E[YD\mid Z=0]}{p_A},\qquad
 \mu_{N0}=\frac{\E[Y(1-D)\mid Z=1]}{p_N},
\]
\[
 \mu_{C1}=\frac{\E[YD\mid Z=1]-\E[YD\mid Z=0]}{p_C},\quad
 \mu_{C0}=\frac{\E[Y(1-D)\mid Z=0]-\E[Y(1-D)\mid Z=1]}{p_C}.
\]
These follow by subtracting the common always-taker or never-taker mixture from each observed arm. Terms for zero-mass groups are interpreted through their mass-weighted numerators, without division by zero.

\begin{proposition}
In the bounded-outcome monotone-instrument class above, with no restrictions on the unobserved counterfactual outcomes of always- and never-takers, the sharp population-effect interval is
\[
 \left[p_C\tau_C+p_A(\mu_{A1}-1)-p_N\mu_{N0},\quad
 p_C\tau_C+p_A\mu_{A1}+p_N(1-\mu_{N0})\right].
\]
\end{proposition}
\begin{proof}
Decompose the population effect as
\[
 p_C\tau_C+p_A(\mu_{A1}-\mu_{A0})
             +p_N(\mu_{N1}-\mu_{N0}).
\]
Only $\mu_{A0}$ and $\mu_{N1}$ are unknown, and each lies in $[0,1]$. The lower endpoint sets these two means to one and zero, respectively; the upper endpoint reverses them. These assignments are compatible with the complete observed law: append deterministic unobserved counterfactuals to the existing type-specific observed potential outcomes and keep assignment independent. Intermediate choices or mixtures attain every value between the endpoints. This proves sharpness, not merely validity of separate coordinate bounds.
\end{proof}

For example, take $(p_A,p_C,p_N)=(0.2,0.3,0.5)$, $\mu_{A1}=0.8$, $\mu_{N0}=0.4$ and $\tau_C=0.2$. The interval is $[-0.18,0.52]$. A positive complier effect therefore does not even sign the population effect in this example. These are hypothetical compatible parameters, not estimates for any election.

\section*{Exclusion sensitivity and weak first stages}
Public encouragement might itself signal endorsement or expected viability. To make the sensitivity algebra exact, allow potential outcomes $Y(d,z)$ and define
\[
 \Gamma=\E[Y(D(1),1)-Y(D(1),0)].
\]
Adding and subtracting $Y(D(1),0)$ yields
\[
 \Delta_Y^Z=\Gamma+p_C\E[Y(1,0)-Y(0,0)\mid C].
\]
If a substantive sensitivity restriction gives $|\Gamma|\leq\delta$, the local spending effect at $z=0$ lies in
$[(\Delta_Y^Z-\delta)/p_C,(\Delta_Y^Z+\delta)/p_C]\cap[-1,1]$.
This is a valid sensitivity interval, not asserted sharp under every additional observation restriction. Its width before clipping is $2\delta/p_C$: a weak spending response magnifies even small direct encouragement effects. More data can shrink sampling error, but cannot remove this stipulated exclusion uncertainty.

\section*{Why controlling for the winner does not separate channels}
Let spending assignment $A$ and latent jurisdiction quality $U$ be independent fair Bernoulli variables. Set winner indicator $V=1\{A=1\text{ or }U=1\}$ and policy $Y=U$, which is unaffected by spending. Among jurisdictions where $V=1$, the mean policy is one when $A=0$ and one half when $A=1$. Conditioning on the winner creates an apparent effect of $-1/2$ although the true total effect is zero.

Even observing randomized spending, the winner and policy does not automatically identify a controlled direct effect. Suppose $V(a)=a$ and the observed policy is $Y=A$. Model P sets $Y(a,v)=a$; model S sets $Y(a,v)=v$. Both yield exactly the same randomized observed law and total effect one. Holding $v$ fixed gives direct effect one in P and zero in S. The data never reveal the off-path combinations needed to distinguish them. A winner-fixing intervention must itself be defensible; this counterexample does not claim that literally appointing the winner is an admissible electoral experiment.

A post-election lobbying experiment among a predetermined set of winners identifies a subsequent lobbying effect for that cohort. It does not retrospectively identify the election-selection channel of earlier campaign spending. Combining the experiments requires a transport and timing model, not merely a regression adjustment.

\section*{Verification and final status}
The bounds preserve the difference between package assignment and realized spending, include zero-mass conventions, and construct compatible endpoint models. Winner conditioning was checked by enumerating the four equally likely $(A,U)$ states. The complete research target still requires a real design, outcome measurements, equilibrium spillover restrictions and an empirically justified law class. The effect of money in politics has not been determined from these benchmark calculations. \textbf{Status: unresolved.}

\section*{Primary references checked}
John M. de Figueiredo and Brian Kelleher Richter (2014), \emph{Advancing the Empirical Research on Lobbying}, Annual Review of Political Science 17, 163--185. Publisher abstract inspected for its identification and measurement agenda. \url{https://www.annualreviews.org/doi/10.1146/annurev-polisci-100711-135308}.

Stephen Ansolabehere, John M. de Figueiredo and James M. Snyder Jr. (2003), \emph{Why is There so Little Money in U.S. Politics?}, Journal of Economic Perspectives 17(1), 105--130. Publisher bibliographic record inspected; no claim of inspecting the full empirical review. \url{https://www.aeaweb.org/articles?id=10.1257/089533003321164976}.

\end{document}
